% Created 2019-03-19 Tue 09:20
\documentclass[11pt]{article}
\usepackage[utf8]{inputenc}
\usepackage[T1]{fontenc}
\usepackage{fixltx2e}
\usepackage{graphicx}
\usepackage{longtable}
\usepackage{float}
\usepackage{wrapfig}
\usepackage{rotating}
\usepackage[normalem]{ulem}
\usepackage{amsmath}
\usepackage{textcomp}
\usepackage{marvosym}
\usepackage{wasysym}
\usepackage{amssymb}
\usepackage{hyperref}
\tolerance=1000
\usepackage{csquotes}
\usepackage{graphicx}
\usepackage{amsmath}
\usepackage{mathtools}
\usepackage{algpseudocode}
\usepackage[margin=0.75in]{geometry}
\usepackage{nth}
\usepackage{tikz}
\usetikzlibrary{positioning,chains,fit,shapes,shapes.misc,calc,arrows,patterns,decorations.pathreplacing}
\usepackage{ccicons}
\usepackage{calc}
\usepackage{fancyhdr}
\usepackage{fancyvrb}
\usepackage{hyperref}
\usepackage{nopageno}
\usepackage{enumitem}
\usepackage{multicol}
\usepackage{braket}
\algnewcommand\algorithmicforeach{\textbf{for each}}
\algdef{S}[FOR]{ForEach}[1]{\algorithmicforeach\ #1\ \algorithmicdo}
\setlength{\columnseprule}{1pt}
\newcommand{\bonusmarksamt}[1]{\textbf{[#1~BONUS~marks]}}
\newcommand{\bonusmarkamt}[0]{\textbf{[1~BONUS~mark]}}
\newcommand{\marksamt}[1]{\textbf{[#1~marks]}}
\newcommand{\markamt}[0]{\textbf{[1~mark]}}
\newcommand{\fillinMC}[1]{\fillinMCmath{\mbox{#1}}}
\newcommand{\fillinMCmath}[1]{\begin{tikzpicture}\draw circle [radius=0.5em];\end{tikzpicture}\ #1}
\newcommand{\fillinMCsoln}[1]{\fillinMCmathsoln{\mbox{#1}}}
\newcommand{\fillinMCmathsoln}[1]{\begin{tikzpicture}\draw[black, fill=blue] circle [radius=0.5em];\end{tikzpicture}\ #1}
\newcommand{\fillinMCAll}[1]{\fillinMCAllmath{\mbox{#1}}}
\newcommand{\fillinMCAllmath}[1]{\begin{tikzpicture}\draw (0,0) rectangle (1em,1em);\end{tikzpicture}\ #1}
\newcommand{\fillinMCAllsoln}[1]{\fillinMCAllmathsoln{\mbox{#1}}}
\newcommand{\fillinMCAllmathsoln}[1]{\begin{tikzpicture}\draw[black, fill=blue] (0,0) rectangle (1em,1em);\end{tikzpicture}\ #1}
\newcommand{\fillinblank}[1]{\fillinblankmath{\mbox{#1}}}
\newcommand{\fillinblankmath}[1]{\begingroup\setlength{\fboxsep}{1em}\setlength{\fboxrule}{2pt}\fbox{\LARGE\phantom{#1}}\endgroup}
\newcommand{\fillinblanksoln}[1]{\fillinblankmathsoln{\mbox{#1}}}
\newcommand{\fillinblankmathsoln}[1]{\begingroup\setlength{\fboxsep}{1em}\setlength{\fboxrule}{2pt}\fbox{{#1}}\endgroup}
\newcommand{\answerbox}[1][3\baselineskip]{%
\noindent\setlength{\fboxrule}{2pt}\framebox[0.7\linewidth]{%
\raisebox{2pt}[0pt][#1]{}%
}\par\medskip%
}
\usepackage{collcell,array}
\newcommand{\shiftdown}[1]{\smash{\raisebox{-.5\normalbaselineskip}{#1}}}
\newcolumntype{C}{>{\collectcell\shiftdown}c<{\endcollectcell}}
\newcommand{\shiftup}[1]{\smash{\raisebox{.5\normalbaselineskip}{#1}}}
%%% Adding Colour to Questions and Answers                                                            
\usepackage{color}
%Questions                                                                                            
\def\question#1{\red{#1}}
% Answers                                                                                             
%Answer in short form                                                                                 
\def\ans#1{\par\blu{Answer: #1}}%With header of Answer                                                
\def\soln#1{\par\blu{#1}} %Without header of Answer                                                   
% Colors                                                                                              
\definecolor{blu}{rgb}{0,0,0.5}
\def\blu#1{{\color{blu}#1}}
\definecolor{gre}{rgb}{0,.3,0}
\def\gre#1{{\color{gre}#1}}
\definecolor{red}{rgb}{0.5,0.0,0}
\def\red#1{{\color{red}#1}}
\def\norm#1{\|#1\|}
%%% End for Colours            
\def\sequence#1#2{\dsequence{#1}{1}{#2}}
\def\dsequence#1#2#3{#1_{#2}, \ldots, #1_{#3}}
\newtheorem{fact}{Fact}
\def\imax{$\texttt{S}_{\max}$}
\def\iprevmax{previous\_$\texttt{S}_{\max}$}
\def\deleted{\texttt{covered}}
%\fancyhead{}\fancyfoot[C]{\tiny This work is licensed under a \href{http://creativecommons.org/licenses/by/4.0/}{Creative Commons Attribution 4.0 International License. \ccby}\\ For license purposes, the author is the University of British Columbia.}\pagestyle{fancyplain}
\makeatletter
\def\@myproof#1[#2]{\trivlist \item [\hskip\labelsep{\bf #1 {#2}: }]}
\def\@begproof#1{\@ifnextchar[{\@myproof{#1}}{\@myproof{#1}[\unskip$\,$]}}
\newenvironment{solution}{\vskip -10pt\leavevmode\@begproof{Solution}}{\endtrivlist}
\def\sopt{S_{\textit{\scriptsize opt}}}
\makeatother
\fancyhead{}\fancyfoot[C]{\tiny This work is licensed under a \href{http://creativecommons.org/licenses/by-nc/4.0/}{Creative Commons Attribution-NonCommercial 4.0 International License. \ccbync}\\ For license purposes, the author is the University of British Columbia.}\pagestyle{fancyplain}
\usepackage{mdframed}
\BeforeBeginEnvironment{minted}{\begin{mdframed}}
\AfterEndEnvironment{minted}{\end{mdframed}}
\date{}
\title{CPSC 320 2018W2: Assignment 5 (Last assignment!)}
\hypersetup{
  pdfkeywords={},
  pdfsubject={},
  pdfcreator={Emacs 25.2.1 (Org mode 8.2.10)}}
\begin{document}

\maketitle
\vspace{-.7in}

Please follow the guidelines given in Assignment 1 for submission to
Gradescope, and for group collaboration. Remember to provide short
justifications for your answers. Submit by the deadline \textbf{Monday 
April 1, 2019 at 10PM}.

Use \LaTeX to prepare your answers. Easiest will be to use the .tex
file provided.  Please enclose each paragraph of your solution in
\verb~\soln{Your solution here...}~.  \soln{Your solution will then appear
in dark blue}, making it a lot easier for TAs to find the parts that
you wrote.

\section{Subset sums mod $n$}
\label{sec-1}
You've now seen examples of dynamic programming problems where the
recurrence describes a number (e.g., the length of a longest common
subsequence), and one example where the recurrence value is Boolean (true or false, indicating
whether or not a there is a subset of a set items whose sum is a given value).
Here you'll work with a problem where the recurrence describes a set.

Let $S = \{x_1, x_2, \ldots, x_n\}$ be a set of nonnegative integers.
We say that a value $v$, with $0 \le v \le n-1$, is \emph{feasible} with respect
to $S$ if for some non-empty subset $R$ of $S$, 
\[ 
\sum_{x \in R} x = v \pmod n.  
\]

\begin{enumerate}
\item For $1 \le i \le n$, let $V(i)$ be the set of feasible values that can
be obtained from $\{x_1, x_2, \ldots, x_i\}$.
\question{Explain why the following recurrence holds for $i \ge 1$, if we define
$V(0) = \emptyset$ (the empty set):}
 \[
    V(i) = V(i-1) + \cup_{v \in V(i-1)} \{ v + x_i \pmod n \} \cup \{ x_i \pmod n \}.
    \]
\soln{
   We can see this recurrence in three parts, the first part $V(i-1)$ ensure that $V(i)$ contains
   the $v_{\text{s}}$ for all subsets of the previous set $\{x_1, x_2 ... x_{i-1}\}$. \\ \\
   With $\cup_{v \in V(i-1)} \{ v + x_i \pmod n \}$ we are addding all the $v_{\text{s}}$ for all the 
   combinations of the previous subsets in $V(i-1)$ with our new value $x_i$, since we only care 
   about the $\pmod n$ of the sum of each subset, we dont need to re calculate the sum over each 
   subset to include $x_i$ we just need the corresponding $v$ of each subset then if we add $x_i$ to 
   it and compute $\pmod n$ it would yield the same result without having to sum over the subset $R$, 
   with this recurrence we are progressively adding to the subsets our new values $x_i$ on each iteration. \\ \\
   Lastly we are missing one subset, we have all the subsets that not include $x_i$ from the 
   first part, and all the subsets that include $x_i$ and one or more $x_j$ in $S$ with $j < i$. 
   But we have not considered the subset $R = \{x_i\}$ and that is what the last part of the 
   recurrence do, $\{ x_i \pmod n \}$, the sum of that set is just $x_i$ so we have include 
   $x_i \pmod n$ into our set $V(i)$. \\
}
\item \question{Design an algorithm that, given a set $S$ of $n$ nonnegative
integers, determines whether 0 is feasible with respect to $S$.} The
algorithm outpus "Yes" or "No".  Your algorithm should run in
$O(n^2)$ time. Your pseudocode can use set options such as $\cup$ as
in the recurrence above.
\soln{
\begin{tabbing}
xxxxxxx\=xxx\=xxx\=xxx\=xxx\=xxxx\= \kill
\> {\bf Algorithm} Subset-Sum-mod-n($S = \{x_1,x_2,\ldots,x_n\}$) \\
\>\> // Returns "Yes" if for some subset $R$ of $S$, $\sum_{x \in R} x = 0 \pmod n$   \\
\>\> Create an array Soln$[0..n]$ \\
\>\> Set Soln$[0] = \emptyset$ \\
\>\> for $i = 1$ to $n$: \\
\>\>\> Soln$[i]$ = \{S$[i] \mod n$\} \\
\>\>\> for $j = 1$ to $|$Soln$[i-1]|$: \\
\>\>\>\> $v$ = Soln$[i-1][j]$ \\
\>\>\>\> Soln$[i]$ = Soln$[i]$ $\cup$ $\{v, v+s[i] \pmod n\}$ \\
\end{tabbing}
}
\item Suppose that indeed there is a subset $R$ of $S$ such that $\sum_{x
   \in R} x = 0 \pmod n$.  \question{Write an algorithm that finds such a subset
$R$}.  You can assume that an array Soln$[0..n]$ has already been
pre-computed, where Soln$[i]$ stores the set $V(i)$ defined in part 1,
and your algorithm can use the array Soln. Your algorithm should
run in time $O(n^2)$.
\end{enumerate}
% JATJ: Complete

\section{NP True or False}
\label{sec-2}
Let $X$ and $X'$ be decision problems, where both problems
have "Yes" instances and "No" instances.
\question{State whether you think each of the following statements must be true,
must be false, or is an open question. Justify your answer.}


\begin{enumerate}
\item \textbf{Statement:}
   If $X \le_p X'$ and $X$ is \emph{not} in NP, then $X'$ is not in NP.
\soln{
   True, if $X'$ is at least as hard as $X$ which is not in NP and also not in P, then it must be that 
   $X'$ has no solution neither in P nor NP. If it was the case that there was a solution for 
   $X'$ that could be checked in polynomial then $X$ would also have a solution that could be checked 
   in polynomial time, but it is not the case since we know that $X$ is not in NP. Hence there is no 
   way to check a solution for $X'$ on polynomial time, therefore $X'$ is not in NP.
   % Open question, we can't confirm that $X'$ is not in NP by just proving that $X'$ which is not in NP is 
   % at least as hard as $X$. This can be clearer if we assign some problem to $X$, like checking if an 
   % array contain a provided element, this would be a linear problem. Knowing that $X'$ is at least as 
   % hard as finding an element in an array, don't provide us with more information about $X'$ it could 
   % easily be infinitely more difficult than P but we don't know.
}
\item \textbf{Statement:}
If $X \le_p X'$, $X$ is in P and $X'$ is in NP, then $X'$ must be in P.
\soln{
   Open question, if I can confirm or deny this, then I would win \$1 million US. By proving this 
   we would prove that P = NP.
}

\item Let
$\overline{X}$ be the \emph{complement} of problem $X$. That is,
instance $I$ of $X$ is a No-instance of $\overline{X}$ if and only
if $I$ is a Yes-instance of $X$. 
For example, the problem $\overline{\mbox{SAT}}$
is the set of \emph{unsatisfiable} 3-Sat Boolean formulas---formulas
for which there is no satisfying truth assignment.
Define $\overline{X'}$ similarly in terms of $X'$. The statement then is:

\textbf{Statement:} If $X \le_p X'$ then $\overline{X} \le_p \overline{X}'$.
\soln{
   False, even if we know that $X \le_p X'$ we can't confirm that proving a not instance is reducible 
   in the same way. So even though $\overline{X}$  has the same context as $X$ they can have complexities, 
   sometimes proving things are no-instances is faster than proving solutions are a yes-instances and 
   vice versa.
}
\end{enumerate}

\section{Common superstrings}
\label{sec-3}
A \emph{common superstring} of the strings $s_1, s_2, \ldots, s_n$ is a
string $T$ such that each $s_i$ is a substring of $T, 1 \le i \le
n$. (Note that the letters of the substring must appear consecutively in
$T$, unlike our definition of a subsequence.) For example, if the
strings are $s_1 = 101, s_2 = 0111$, and $s_3 = 010$ then the strings
010111 and 0111010 are both common superstrings of $s_1, s_2$, and
$s_3$.  

An instance of the \emph{Shortest Common Superstring} problem is a set of
strings $\{s_1, s_2, \ldots, s_n\}$ and a number $K$.
The problem is to determine whether there is a
common superstring of $s_1, s_2, \ldots, s_n$ that has length at most
$k$.  Assume that none of the strings $s_i$ is a substring of another
string $s_j$ (otherwise $s_i$ can be removed from the list of strings
without changing the length of the shortest common superstring). This
problem has been important in the context of developing algorithms for
genome sequence assembly.

An instance of the \emph{Traveling Salesperson Problem} (TSP) is a set $\{
C_1, C_2, \ldots, C_n \}$ of cities, nonnegative integer costs
$c_{ij}$ of traveling from city $C_i$ to city $C_j$ (or from $C_j$ to
$C_i$; the cost is the same either way), and a number $K$.  The
problem is to determine whether there is a tour of the cities of cost
at most $K$.  A \emph{tour} is a path that starts at some city, say $C_1$,
travels through every other city exactly once, and returns to the
starting city $C_1$. The tour cannot go through a city more than once
(except for $C_1$ which appears both at the start and at the end of
the tour).

\begin{enumerate}
\item \question{Describe a polynomial time reduction from the Shortest Common
Superstring problem to the traveling salesperson problem (TSP).}
\soln{
   We can set every string $s_i$ to be a "city" $c_i$, $1 \le i \le n$.
   Then we create a $n$ $c_{ij}$ costs, each of this costs will be 
   the minimum number of characters needed to create a super string 
   of the strings $\{s_i, s_j\}$, i.e. for the $s_i = 010$, $s_i = 0111$ 
   the cost $c_{ij}$ would be 2, because we can use the 0 on both on the 
   left side of the $s_j$ and just add the 2 remaining characters $01$ on 
   $s_i$. \\ \\
   We can create the each city $c_i$ in linear time, just by iterating 
   over the strings. After the creation of our costs, would take us $n^2$
   time because we need to calculate a $n-1$ costs for each string $s_i$.
}

\item \question{Is the Shortest Common Superstring problem in NP? Justify your answer.}
\soln{
   Yes, we reduced Shortest Common Superstring in polynomial time to 
   Traveling Salesperson Problem, and we know that TSP can be verified 
   in polynomial time. Which means that we can verify SCS in polynomial time 
   by using TSP to verify it. Therefore SCS is in NP.
}

\item \question{Does the reduction of part 1, plus your answer to part 2, imply
that the Shortest Common Superstring problem is NP-complete?  Why
or why not?}
\soln{
   No, even we know for our answer in part 2 that SCS is in NP, we don't 
   have enough facts to prove that SCS is in NP-Complete. What we have 
   proved is that TSP is at least as hard as SCS, but to prove that 
   SCS we need to prove that SCS is at least as hard as any NP-Complete 
   problem, we could achieve this by reversing the direction of the reduction 
   reducing TSP to SCS, then we will know that SCS is at least as hard as TSP 
   which is an NP-Complete problem, hence SCS must be an NP-Complete.
}
\end{enumerate}

\section{Maximum Cuts (from Coulter Beeson)}
\label{sec-4}
In this problem you'll define and show that the MaxCut problem is
NP-complete, using a reduction from a known NP-Complete problem,
Partition.  We'll first define Partition, which is already familiar to you, 
and then move on to MaxCut.

\textbf{Partition} is another name for the decision problem we saw in class,
where thieves try to evenly share their loot. An instance of this
problem is a set $I = \{v_1, v_2 ... v_n\}$ of values for $n$ items.
$I$ is a Yes-instance if the items can be fairly split, that is, if we
can partition the set $[1..n]$ into two nonempty subsets $S$ and
$\overline{S}$ such that $\sum_{i\in S}v_i = \sum_{j \in \overline{S}}
v_j$.

A \textbf{cut} in an undirected graph $G=([1..n],E)$ divides the set of nodes $[1..n]$ into
two non-empty disjoint subsets. That is, a cut is a pair $(S,
\overline{S})$ where $\overline{S} = V-S$
and $0 < |S| < n$. If weight $w(i,j)$ is associated with each edge $(i,j)$ of the graph,
then the weight of the cut $(S,\overline{S})$ is the sum of
the weights of edges that cross the cut. That is,
\[
w(S,\overline{S}) = \sum_{i \in S, j \in \overline{S}} w(i,j).
\]
We would like to show that there is there is unlikely to be an
efficient algorithm that takes as input a graph $G$ and weight
function $w$, and finds a cut of maximum weight. However
NP-Completeness is only defined for decision problems.

\begin{enumerate}
\item \question{Write down a decision version of the problem, which
we'll call MaxCut. What is an instance $I$ of MaxCut? Which are the
Yes-instances of MaxCut?}
\soln{
   The decision would be, is there a cut with a weight of at least $k$. Then an instance $I$ 
   would be the graph $G=([1..n],E)$, a set of weights $W$ for each edge in $G$, and a $k$ 
   that would represent the boundary to determine if $I$ is a yes-instance or a no-instance.
}

\item \question{Show that $\text{MaxCut} \in \text{NP}$.}
\soln{
   Given a certificate $C$ which is a set of edges in the MaxCut for a yes-instance $I$, 
   we can verify that is truthly a yes-instance by: 
   \begin{enumerate}
      \item Verifying that every edge $(i,j)$ in $C$ is in MaxCut, and there are no duplicates on $C$
      \item Sum each $w(i,j)$ for every edge $(i,j)$ in $C$. If this sum is at least $k$ then 
      $I$ is truthly a yes-instance
   \end{enumerate}
   This can be done in polynomial time, because the edges in $C$ are upper bounded by $n^2$, because 
   there can't be duplicates on $C$. Hence MaxCut is in NP.
}

\item Here you'll reason about a useful building block for a reduction
from Partition to MaxCut.
Let $I = \{v_1,v_2,\ldots,v_n\}$ be an instance of Partition, with
$\sum_{i=1}^n v_i = 2V$.  From $I$ we'll build a graph $G =
   ([1..n], E)$ where $E$ has an undirected edge between every pair of
nodes $(i,j)$, with weight $w(i,j) = v_i v_j$.
\question{Show that if $(S, \overline{S})$ is a partition of $[1..n]$ such that
$\sum_{i\in S} v_i = V$, then in $G$, the weight of the cut $(S,S')$
is $V^2$.}
\soln{
   Consider an $I$ with $n=4$ with elements $\{v_1, v_2, v_3, v_4 \}$, arbitrary assuming that a subset $S$ has 
   the elements $\{ v_1, v_2 \}$, therefore $\overline{S} = \{ v_3, v_4\}$. Since we are setting the 
   weights $w(i,j)$ to be a multiplication between $v_i$ and $v_j$ on $G$, we will have that the sum of 
   the weights in the cut is, $v_1v_3 + v_1v_4 + v_2v_3 + v_2v_4$, we can think about this as a 
   distribution of two factors, the sets $S, \overline{S}$. So we end up with $(v_1+v_2) * (v_3+v_4)$, which 
   still is equivalent to the weight of the cut. \\ \\
   But we can go further with this, we know that $v_1+v_2 = \sum_{i \in S}^n v_i$ and
   $v_3+v_4 = \sum_{i \in \overline{S}}^n v_i$, so we can reduce the weight of the cut to 
   $(\sum_{i \in S}^n v_i) * (\sum_{i \in \overline{S}}^n v_i)$, furthermore we are given that 
   $\sum_{i\in S} v_i = V$ which yields, weight of cut = $V * (\sum_{i \in \overline{S}}^n v_i)$. Lastly 
   we know that the sum of all elements in $I$ add up to $2V$, since $\overline{S}$ is the complement 
   of $S$, $2V = V + \sum_{i \in \overline{S}}^n v_i$, therefore $ \sum_{i \in \overline{S}}^n v_i = V$
   which give us the weight of the cut $ = V*V = V^2$.
}

\item \question{Show that if $(S, \overline{S})$ is a partition of $[1..n]$ such that
$\sum_{i\in S} v_i \neq V$, then in the graph $G$ of part 3, the weight of the cut $(S,\overline{S})$
is less than $V^2$.}
\emph{Hint:} use the fact that if $x + y = 2V$ then $x y$ attains its unique maximum
when $x = y = V$.
\soln{
   Lets consider a similar scenario like the above one, $S = \{ v_1, v_2 \}$ and $\overline{S} = \{ v_3, v_4\}$.
   Now we are given that $\sum_{i\in S} v_i \neq V$, so lets pick an arbitrary number $c$, $c>0$ to be the 
   difference between $\sum_{i\in S} v_i$ and $V$, so we have $\sum_{i\in S} v_i = V - c$ and 
   $\sum_{i\in \overline{S}} v_i = V + c$ because the difference must be added to the complement for 
   consistency. \\ \\
   Given these expressions we can follow the procedure of the previous question to arrive to the 
   conclusion that the weight of the cut $= (V-c)*(V+c)$, given that the $I$ has $\sum_{i=1}^n v_i = 2V$ 
   and the fact that if $x + y = 2V$, the unique maximum is when $x = y = V$. Then the weight of the 
   cut  $(V-c)*(V+c) < V^2$. Because $(V-c) \neq (V+c) \neq V$ for $c > 0$, which has to be greater 
   than 0 to meet the initial assumtion $\sum_{i\in S} v_i \neq V$.
}

\item \question{Describe an efficient reduction that maps instances $I$ of Partition to instances $I'$ 
of MaxCut.}
You can describe your reduction in words or implement it in pseudocode.
\soln{
   To reduce Partition to MaxCut, we can create a graph $G$ containing a node $a_i$ for each element in the Partition set 
   $\{v_1,v_2,...v_n\}$, $1 \leq i \leq n$. Then create an edge between all $a_i$ and $a_j$ nodes 
   and set $w(i,j) = v_i v_j$, $1 \leq i \leq n$, $1 \leq j \leq n$, $i \neq j$. Lastly set 
   $k = ( \frac{1}{2} \sum_{i=1}^n v_i)^2$.
}

\item \question{Explain why your reduction runs in polynomial time.}
\soln{
   This reduction takes linear time to create the nodes, quadratic time for creating 
   the edges and weights between every $x_i$ and $x_j$, and again linear time to calculate $k$. 
   The total running time is $n^2$ which is polynomial time. 
}

\item \question{Show that if $I$ is a Yes-instance of Partition then $I'$ is a
Yes-Instance of MaxCut.} (Part 3 should be useful here, as well as
your definition of part 1.)
\soln{
   If $I$ is a yes-instance then the sum of each $v_i$ in $I$ will be divisible by 2, because we need 
   to separate $I$ in two equal parts to be a yes-instance, this implies that $\sum_{i}^n v_i = 2V$ 
   hence  $\sum_{i\in S} v_i = \sum_{i\in \overline{S}} v_i = V$. \\ \\
   From part 3 we know that if the sum of the subset $S$ equal $V$ then the weight of the maximum 
   cut would be $V^2$, so now we just need to compare $k$ and $V^2$, so we have the following 
   comparison $V^2 = ( \frac{1}{2} \sum_{i=1}^n v_i)^2$, which if we replace the sum of every $v_i$ in $I$
   with $2V$ yields $V^2 = ( \frac{1}{2} (2V))^2 \rightarrow V^2 = (V)^2$, which is true. Therefore 
   $I'$ is also a yes-instance.
}

\item \question{Show that if $I'$ is a Yes-instance of MaxCut then $I$ is a Yes-Instance of Partition.}
(Part 4 should be useful here.)
\soln{
   If $I'$ is a yes-instance then the weight of the max cut $ = ( \frac{1}{2} \sum_{i=1}^n v_i)^2$, 
   which can also be represented as MaxCut $ = V^2$. We know from part 4 that there is no subset 
   that for any subset $S$ that add up to some number different than $V$, the weight of MaxCut would 
   be strictly less than $V^2$. But in this case we know that it is $ = V^2$, therefore $I$ must 
   also be a yes-instance.
}

\item In class we saw a dynamic programming algorithm for the Partition
problem, with running time $\Theta(nV)$, where $2V = \sum_{i=1}^n v_i$.
\question{Why does this not imply that $\text{P}=\text{NP}$?} (Think what is
the size of a Partition instance; is the running time of the dynamic
programming algorithm polynomial or exponential in the input size?)
\soln{
   This do not imply that P = NP, the running time of that algorithm grows not only with the size of 
   the input, but with the size of the data in the input. So we can have an instance $I$ with $n=5$ 
   and another $\hat{I}$ with $n=5$, but they would have different running times based on the value 
   of $V_I$ and $V_{\hat{I}}$.
}

\item \textbf{Bonus} \question{Prove the hint given in part 4.}
\soln{
   To prove that $xy$ reach their unique maximum if $x+y=2V$ when $x=y=V$. Consider the following 
   expression to prove, $(x-c)(y+c) < xy$ for some $c \neq 0$, $c < V$, if we distribute that expression 
   we get: 
   \begin{center}
      \[(x-c)(y+c) < xy\]
      \[xy + xc - yc - c^2 < xy\]
      \[xc - yc - c^2 < 0\]
      \[x - y - c < 0\]
      \[x - c < y\]
      By considering $x+y=2V$
      \[x - c < 2V -x\]
      \[2x - c < 2V \]
      We know for $x+y=2V$, that $x$ is at most $2V$, so by replacing it
      \[2V - c < 2V \]
      Which yields true, due to our constraints on $c$, $c \neq 0$, $c < V$. Therefore, our initial 
      clause $(x-c)(y+c) < xy$ is also true. Hence $xy$ reach the maximum when $x = y = V$, when 
      $c=0$.
   \end{center}
}
\end{enumerate}
\end{document}
